% --- Archon physics formalization source begin ---
% archon:physics
% archon:chemistry
% archon:covers IChO2026Problems/problem_icho_2026_t9_a9.lean
% archon:source-report reports/icho_2026/problem_icho_2026_t9_a9.source.json
% archon:problem-id icho_2026_t9
% archon:part-id T9-A9

\section{Icho Chemistry problem icho\_2026\_t9\_a9}

\paragraph{Problem source.}
\#\# Source
58th International Chemistry Olympiad, Tashkent, Uzbekistan, 2026.
Official-materials index: https://www.icho2026.uz/problems
English problem PDF: ../icho\_2026\_source/raw/theory\_problem.pdf
English solution/rubric PDF: ../icho\_2026\_source/raw/theory\_solution.pdf
Problem PDF page 88; printed local question page 5; solution starts on PDF page 90.
The page PNG is primary visual evidence for formulas, structures, tables, and figures.

\#\# T9. Cyclodextrin Chemistry
T9. Cyclodextrin Chemistry 7\% Cyclodextrins (CD) are a family of cyclic oligosaccharides, consisting of glucose subunits joined by α1,4-glycosidic bonds. The three most common cyclodextrins (α-, β-, γ-cyclodextrin) contain 6, 7, and 8 α-D-glucopyranoside units, respectively. 9.1 Calculate the molar mass of β-CD. Assume M𝑤(glucose) = 180.16 g mol−1. CDs combine macrocyclic properties with glucose chemistry, which was demonstrated in the synthesis of X. Primary and secondary hydroxy groups show different reactivity due to being a part of CDs. Stoddart et al. synthesised compound K, where only one OH group remained free per glucopyranoside unit. equiv.= equivalent 9.2 Tick the favorable chair conformation and draw the structure of K by completing the CD template. Show stereochemistry unambiguously. 9.3 Determine the ring size, 𝑟𝑠, of macrocycle X. Give the number of stereocentres, 𝑠𝑐, in X. α-Cycloaltrin, a synthetic analogue of α-CD, was synthesised as follows. 9.4 Draw the structure of Y with stereochemistry by completing the CD template. In 2000, Sinay et al. demonstrated that a single protic group (NH and OH) at unit 1 directs the next reductive debenzylation of a primary OH group to unit 4 in the macrocyclic ring (or to unit 3 if the unit 4 position is not available). This allowed the synthesis of the following β-CD dimer as a mixture of constitutional linkage isomers. 9.5 Draw the structure of L on the β-CD template. Fill in all the boxes. 9.6 Determine the number of isomers of the β-CD dimer that can form during the synthesis by Sinay et. al. To confirm the positions of debenzylation, L was subjected to the “hexo-5-enose degradation” reaction, as shown below. The products were then analysed by mass spectrometry. 9.7 Calculate the 𝑚/𝑧value for the two [M+Na]+ peaks that were observed for the degradation products from L. Use integer values of atomic mass. Later, Sollogoub et al. applied the ability of alkenes to direct the debenzylation to the “ortho”- glucopyranose unit for the synthesis of hexadifferentiated α-CD S: 9.8 Draw the structures of O-S on the α-CD templates. Fill in all the boxes. You may draw structures outside the boxes if appropriate. Assume 1,2-unit modification happens clockwise, while 1,3-unit modification happens counterclockwise. Protic groups have stronger directing effects than alkenes.

\#\# Current subquestion 9.9
Determine the number of all possible arrangements of the functional groups on a hexadifferentiated α-CD (assume only the CH2OH groups have been modified).

\#\# Dataset metadata
IChO 2026; T9; raw rubric points=4; kind=theory; formalization\_ready=true.

\paragraph{Shared problem context.}
T9. Cyclodextrin Chemistry 7\% Cyclodextrins (CD) are a family of cyclic oligosaccharides, consisting of glucose subunits joined by α1,4-glycosidic bonds. The three most common cyclodextrins (α-, β-, γ-cyclodextrin) contain 6, 7, and 8 α-D-glucopyranoside units, respectively. 9.1 Calculate the molar mass of β-CD. Assume M𝑤(glucose) = 180.16 g mol−1. CDs combine macrocyclic properties with glucose chemistry, which was demonstrated in the synthesis of X. Primary and secondary hydroxy groups show different reactivity due to being a part of CDs. Stoddart et al. synthesised compound K, where only one OH group remained free per glucopyranoside unit. equiv.= equivalent 9.2 Tick the favorable chair conformation and draw the structure of K by completing the CD template. Show stereochemistry unambiguously. 9.3 Determine the ring size, 𝑟𝑠, of macrocycle X. Give the number of stereocentres, 𝑠𝑐, in X. α-Cycloaltrin, a synthetic analogue of α-CD, was synthesised as follows. 9.4 Draw the structure of Y with stereochemistry by completing the CD template. In 2000, Sinay et al. demonstrated that a single protic group (NH and OH) at unit 1 directs the next reductive debenzylation of a primary OH group to unit 4 in the macrocyclic ring (or to unit 3 if the unit 4 position is not available). This allowed the synthesis of the following β-CD dimer as a mixture of constitutional linkage isomers. 9.5 Draw the structure of L on the β-CD template. Fill in all the boxes. 9.6 Determine the number of isomers of the β-CD dimer that can form during the synthesis by Sinay et. al. To confirm the positions of debenzylation, L was subjected to the “hexo-5-enose degradation” reaction, as shown below. The products were then analysed by mass spectrometry. 9.7 Calculate the 𝑚/𝑧value for the two [M+Na]+ peaks that were observed for the degradation products from L. Use integer values of atomic mass. Later, Sollogoub et al. applied the ability of alkenes to direct the debenzylation to the “ortho”- glucopyranose unit for the synthesis of hexadifferentiated α-CD S: 9.8 Draw the structures of O-S on the α-CD templates. Fill in all the boxes. You may draw structures outside the boxes if appropriate. Assume 1,2-unit modification happens clockwise, while 1,3-unit modification happens counterclockwise. Protic groups have stronger directing effects than alkenes.

\paragraph{Current subquestion.}
Determine the number of all possible arrangements of the functional groups on a hexadifferentiated α-CD (assume only the CH2OH groups have been modified).

\paragraph{Recorded answer/context.}
If we have 6 linear positions for 6 coloured balls, the total number of arrangements will be 6! = 720. But in α-CD, positions are in a circle, which makes all arrangements below equal. The correct answer 6!/6 = 120. Correct answer 4 points; if the answer 720 or 60– 2 points, if answer 320 - 1 point. All other answers - 0. 4.0 pt

\paragraph{Figure/image paths.}
\begin{itemize}
\item ../icho\_2026\_source/image/T9\_page-5.png
\item ../icho\_2026\_source/image/T9\_page-4.png
\end{itemize}

\paragraph{Formalization target.}
create a compiling Lean file with sorry bodies at `IChO2026Problems/problem\_icho\_2026\_t9\_a9.lean`.
The Lean declarations must preserve every mathematical object, quantifier, hypothesis, side condition, approximation/error guarantee, and requested conclusion from the source statement.
Use Mathlib/Physlib/Chemistry names found through LeanExplore where available. Prefer the configured benchmark Base library; introduce a local definition only when the source genuinely requires an object that the environment does not expose.

\begin{theorem}\leanok
[Icho Chemistry formalization target]
\label{thm:physics:icho_2026_t9_a9:target}
\lean{IChO2026Problems.T9A9.number_of_hexadifferentiated_alphaCD_arrangements}
The linear placement count and alpha-cyclodextrin symmetry determine the number
of hexadifferentiated arrangements.
\end{theorem}
\begin{proof}

\leanok
Assigning six distinct functional groups to six labelled primary
CH$_2$OH sites is a bijection, so there are $6!=720$ linear placements.  The
six rotations of the cyclic \(\alpha\)-CD framework act on these placements.
No nonidentity rotation can fix a bijection, since it moves a site whose
distinct label would then have to occur at two sites.  Every rotation orbit
therefore has six placements.  The source identifies rotations only, not
reflections, so the number of physical arrangements is $720/6=120$.
\end{proof}
% --- Archon physics formalization source end ---
